\section{Integrating the unitary modules}
\label{sec:integration}

We now construct the local algebra actions on a unitary module.
The input is the bounded regularization of a charged field from
Lemma~\ref{lem:charged-tools}.  We first show that this regularization
intertwines the observables in a disjoint interval as closed operators.
Creation fields will then identify the observable algebra on the module
with the vacuum algebra.  Only the already given positive form on the
module is used in this construction.

For unitary modules $N,K$ and an interval $J$, define the diagonal
observable algebra
\begin{equation}
 \mathcal B_{N,K}(J)=W^*\!\left\{
 \overline{Y^N(v,g)}\oplus\overline{Y^K(v,g)}:
 v\in V\text{ homogeneous},\ 
 g\in C_c^\infty(J)\right\}.
 \label{eq:diagonal-observable-algebra}
\end{equation}
As for $A_V$, an unbounded generator means its closed graph data,
or equivalently its bounded transform and polar data.  If
$T:\mathcal H_N\to\mathcal H_K$ is a closed densely defined operator,
its closed lower-corner operator on $\mathcal H_N\oplus\mathcal H_K$ is
\[
 \widetilde T=\begin{pmatrix}0&0\\T&0\end{pmatrix},
 \qquad\Dom\widetilde T=\Dom T\oplus\mathcal H_K.
\]
We say that $T$ \emph{strongly intertwines the observables in $J$}
when $\widetilde T$ is affiliated with $\mathcal B_{N,K}(J)'$.

\begin{proposition}[Strong intertwining from annulus regularization]
\label{prop:strong-intertwining}
Let $V$ be simple, unitary and strongly rational, and let $M,N,K$ be
grading-restricted unitary $V$-modules, with $N$ strongly integrable.
Let $\Phi$ be an intertwining operator of type $\binom{K}{M\,N}$,
and let $a\in M$ be homogeneous.  Assume that $\Phi(a,z)$ satisfies
the transport domination hypothesis of Lemma~\ref{lem:charged-tools}.
Smear $\Phi(a,z)$ against a smooth test $f$ supported in an interval $I$, with
a fixed argument branch.  If $J$ is disjoint from $I$, then the closed
smeared operator $\overline{\Phi(a,f)}$ strongly intertwines the
diagonal smeared observables of $N$ and $K$ supported in $J$.
This holds for every homogeneous observable in $V$ and for every such
choice of tests and argument branches.
\end{proposition}

\begin{proof}
Write $\mathcal D_N=\mathcal H_N^\infty$ and
$\mathcal D_K=\mathcal H_K^\infty$.
The polynomial mode estimates imply that smeared fields and their formal
adjoints map the corresponding smooth spaces continuously to smooth
spaces.  Spectral truncations converge in all the smooth seminorms.
Consequently a smear has the same closure whether it is initially
defined on finite-energy vectors or on smooth vectors.  These facts
apply both to the observables and, by Lemma~\ref{lem:charged-tools},
to the charged field.

Choose the localized annulus $\underline A$ in that lemma with its dent
containing $\operatorname{supp}f$ and localized in $I$.
Its action $\underline A_N$ is a
bounded injective operator with dense range, and
\begin{equation}
 T_f=\overline{\Phi(a,f)}\underline A_N
 \label{eq:bounded-charge-regularization}
\end{equation}
is bounded on all of $\mathcal H_N$.  Both $T_f$ and $T_f^*$ preserve
the respective smooth spaces.  On $\mathcal D_K$ the adjoint is
$\underline A_N^*\Phi(a,f)^\dagger$, where the dagger denotes the formal smeared
adjoint.  These are assertions about the full bounded regularization,
not merely about its restriction to finite energies.

For homogeneous $v\in V$ and a smooth $g$ supported in $J$, put
\[
 S_N=\overline{Y^N(v,g)},\qquad
 S_K=\overline{Y^K(v,g)}.
\]
On the smooth spaces, the intertwining-operator Jacobi identity gives
\[
 Y^K(v,g)\Phi(a,f)=\Phi(a,f)Y^N(v,g).
\]
To pass from the formal identity to this equality, multiply the two
field orderings by a sufficiently high integral power of the coordinate
difference.  Polynomial energy bounds make the resulting series
smooth-domain distributions.  Since the two test supports are
disjoint, one can divide the product test by that power.  The same
argument applies to the adjoint relation, whose observable test has
the same support.  This is the weak intertwining property in
\cite[Theorem~3.5]{GuiUnbounded}.

The operator $\underline A_N$ and its adjoint commute strongly with these
observables: $\underline A_N$ is localized in $I$, whereas $N$ is already
integrated and the observables are localized in $J$.
Combining this with the weak identities gives
\begin{equation}
 S_KT_f\xi=T_fS_N\xi\quad(\xi\in\mathcal D_N),
 \qquad
 S_NT_f^*\eta=T_f^*S_K\eta\quad(\eta\in\mathcal D_K).
 \label{eq:bounded-charge-core-commutation}
\end{equation}
For the second equality one can pair the first weak intertwining
identity for the \emph{formal adjoint observable}.  More explicitly,
let $S_N^\dagger$ and $S_K^\dagger$ denote its actions on the smooth
spaces.  The invariant adjoint formula expresses it as a finite sum
of smears of PCT-transformed states and their $L_1$-descendants, with
tests still supported in $J$.  Thus the first equality also gives
$T_fS_N^\dagger=S_K^\dagger T_f$ on $\mathcal D_N$.
For $\xi\in\mathcal D_N$ and $\eta\in\mathcal D_K$, using our
linear-first inner product convention, we obtain
\[
\begin{aligned}
 \langle S_NT_f^*\eta,\xi\rangle
 &=\langle T_f^*\eta,S_N^\dagger\xi\rangle
  =\langle\eta,T_fS_N^\dagger\xi\rangle\\
 &=\langle\eta,S_K^\dagger T_f\xi\rangle
  =\langle S_K\eta,T_f\xi\rangle
  =\langle T_f^*S_K\eta,\xi\rangle.
\end{aligned}
\]
All vectors in these pairings are smooth.  This proves the second
equality in \eqref{eq:bounded-charge-core-commutation} without
identifying a formal-adjoint closure with the full Hilbert-space
adjoint.

We spell out the passage to closed operators.  On
$\mathcal H_N\oplus\mathcal H_K$, let
\[
 D=S_N\oplus S_K,\qquad
 Q_f=\begin{pmatrix}0&0\\T_f&0\end{pmatrix}.
\]
Approximation in the graph norms, using the smooth cores, shows that
both $Q_f$ and $Q_f^*$ preserve $\operatorname{Dom}D$ and commute with
$D$ there.  Hence $Q_f\oplus Q_f$ and its adjoint preserve the closed
graph of $D$.  They therefore commute with its orthogonal projection.
The matrix entries of this graph projection determine the usual
bounded transform and polar data of $D$.  Thus $Q_f$ commutes with the
von Neumann algebra generated by $D$.

As $v$ and $g$ vary, the preceding argument proves
$Q_f\in\mathcal B_{N,K}(J)'$.

It remains to remove the annulus.  The first-summand compression of
$\mathcal B_{N,K}(J)$ is $\pi_J^N(A_V(J))$: its normal range is
generated by the module smears, which are the images of the vacuum
generators under strong integration.  Since $\underline A_N$ is
localized in $I$, both $\operatorname{diag}(\underline A_N,1)$ and
its adjoint commute with $\mathcal B_{N,K}(J)$.  The closed inverse
$\underline A_N^{-1}$ is densely defined on
$\operatorname{ran}\underline A_N$, and inversion preserves the
resulting affiliation with the commutant.  Therefore
\[
 Q_f\begin{pmatrix}\underline A_N^{-1}&0\\0&1\end{pmatrix}
\]
is densely defined, and its domain and action commute with every
unitary in $\mathcal B_{N,K}(J)$.  It is closable: by
\eqref{eq:bounded-charge-regularization} it is a restriction of the
closed operator with lower corner $\overline{\Phi(a,f)}$ and all other
corners zero.  Its domain contains the finite-energy core in the first
summand, since $N\subset\operatorname{ran}\underline A_N$ by
Lemma~\ref{lem:annulus-range}.  Its closure is consequently exactly
that closed lower-corner operator.  This proves its affiliation with
$\mathcal B_{N,K}(J)'$.

This affiliation is the strong intertwining assertion, including the
diagram convention used in \cite[Section~3.3]{GuiUnbounded}.
Explicitly, on two copies of $\mathcal H_N\oplus\mathcal H_K$, repeat
the charged corner diagonally and place a diagonal observable in a
matrix corner between the copies.  The latter is affiliated with
$M_2(\mathcal B_{N,K}(J))$, and the repeated charged corner with its
commutant.  Hence the diagram commutes strongly.  All operations above
are valid for each fixed branch on $I$.
\end{proof}

We next use creation fields to show that the vacuum compression of a
diagonal observable algebra loses no information.

\begin{lemma}[Totality of local creation ranges]
\label{lem:creation-totality}
Let $V$ be simple, unitary and strongly rational, let $M$ be a
grading-restricted unitary $V$-module, and let $I$ be a nonempty
interval.  For every homogeneous $a\in M$ and
$f\in C_c^\infty(I)$, form the bounded creation regularization
\[
 T_{a,f}=\overline{\mathcal Y_M^+(a,f)}\underline A_V
\]
of Lemma~\ref{lem:charged-tools}; the localized annulus may depend on
the test.  The closed linear span of their ranges is $\mathcal H_M$.
\end{lemma}

\begin{proof}
The transport domination needed here holds with equality for creation
and vacuum source.  It uses the given form on $M$, and assumes neither
integration of $M$ nor positivity of a general fusion product.
Since $\Omega\in\operatorname{ran}\underline A_V$, each range contains
$\mathcal Y_M^+(a,f)\Omega$.

Suppose $\eta\in\mathcal H_M$ is orthogonal to all these vectors.
For each homogeneous $a$, consider
\[
 F_a(z)=\langle e^{zL_{-1}}a,\eta\rangle,
 \qquad |z|<1.
\]
Creation gives
$\mathcal Y_M^+(a,z)\Omega=e^{zL_{-1}}a$.
The charged energy estimate bounds the Taylor coefficients of $F_a$
by a polynomial in their index.  Its boundary value is therefore a
distribution with only nonnegative Fourier frequencies.  Orthogonality
to all the creation smears says that this distribution vanishes on
$I$.  The smooth weight or branch factor in the smearing convention
can be absorbed into the test on this interval.

For completeness, such a one-sided distribution is zero if it vanishes
on a nonempty arc.  Convolve it with smooth approximate identities
supported in increasingly small arcs.  Each convolution still has
only nonnegative frequencies and vanishes on a smaller nonempty arc.
It is the smooth boundary value of a holomorphic function in the disc.
The Schwarz reflection principle across an arc on which that function
is zero, followed by analytic uniqueness, shows that the convolution
is identically zero.  Passing to the distributional limit proves the
claim.  Thus $F_a=0$, and its constant coefficient gives
$\langle a,\eta\rangle=0$.  Homogeneous vectors span $M$, which is
dense in $\mathcal H_M$; hence $\eta=0$.
\end{proof}

Local-intertwiner criteria for module integration are given in
\cite[Proposition~4.7]{GuiUnitaryII} and
\cite[Lemma~4.27]{GuiExtensions}. The following diagonal-algebra proof
uses the bounded creation operators and their total ranges to establish
the representations together with equality of the full closed graphs.

\begin{theorem}[Strong integration of every unitary module]
\label{thm:integration}
Let $V$ be simple, unitary and strongly rational.  Every
grading-restricted unitary $V$-module $M$ has a unique compatible
family of normal representations
\[
 \pi_J^M:A_V(J)\longrightarrow B(\mathcal H_M)
\]
such that, for every homogeneous $v\in V$ and
$g\in C_c^\infty(J)$,
\begin{equation}
 \pi_J^M\!\left(\overline{Y^V(v,g)}\right)
       =\overline{Y^M(v,g)}.
 \label{eq:exact-cwx-smear}
\end{equation}
Equality here includes the domains of the closed affiliated operators.
The canonical Virasoro action of this net representation is the
integrated module Virasoro action; in particular its rotations are
$e^{itL_0}$ with the module's conformal-energy normalization.
\end{theorem}

\begin{proof}
Module energy bounds follow from Theorem~\ref{thm:energy}, and the
vacuum net exists by Corollary~\ref{cor:locality}.
Fix $J$ and form $\mathcal B(J)=\mathcal B_{V,M}(J)$ as in
\eqref{eq:diagonal-observable-algebra}, including every homogeneous
observable.  Its two Hilbert summands reduce all its generators.
Vacuum restriction is therefore a normal unital homomorphism
\[
 r_0:\mathcal B(J)\longrightarrow A_V(J).
\]
It is onto.  Indeed, compression of the graph data of a direct sum
gives the graph data of its first summand, and these compressed data
generate $A_V(J)$.  The range of a normal homomorphism of von Neumann
algebras is a von Neumann algebra, so it contains that generated
algebra and is equal to it.

The kernel of $r_0$ is $\mathcal B(J)q$ for a central projection $q$.
Choose any interval $I$ disjoint from $J$.  Apply
Proposition~\ref{prop:strong-intertwining} to creation, whose source
is the integrated vacuum.  Its bounded regularizations commute, as
lower corners, with $\mathcal B(J)$.  Since $q$ annihilates the vacuum
summand, its restriction $q_M$ to the other summand satisfies
\[
 q_MT_{a,f}=T_{a,f}q_V=0.
\]
Lemma~\ref{lem:creation-totality} gives $q_M=0$, hence $q=0$.
Thus $r_0$ is a normal isomorphism.  Its inverse is normal as well:
an order isomorphism preserves suprema of bounded increasing positive
nets.  If $r_M$ denotes restriction to the module summand, define
\[
 \pi_J^M=r_Mr_0^{-1}.
\]
For $J_1\subset J_2$ the diagonal algebras are nested, and both
compression maps respect the inclusions.  It follows that
$\pi_{J_2}^M|_{A_V(J_1)}=\pi_{J_1}^M$.

We verify the exact closed-operator identity.
Let $S_V=\overline{Y^V(v,g)}$ and
$S_M=\overline{Y^M(v,g)}$.
The graph projection of $S_V\oplus S_M$ has its matrix entries in
$\mathcal B(J)$.  Compressing its entries to the two summands yields
the graph projections of $S_V$ and $S_M$, respectively.
Consequently the matrix amplification of $\pi_J^M$ sends the entire
graph projection of $S_V$ to that of $S_M$.  This is precisely the
normal image of an affiliated closed operator, and proves
\eqref{eq:exact-cwx-smear}.  In particular the argument identifies
closed graphs, rather than merely operators on a dense common domain.

The graph data of the vacuum smears generate each $A_V(J)$, so the
same identity also proves uniqueness.  With the module energy bounds,
it is the defining property of strong integrability and of the CWX
action; see \cite[Definition~4.22]{GuiExtensions}.
Applying the identity to the conformal vector identifies all localized
stress-energy fields.  Their canonical centrally extended Virasoro
actions agree by \cite[Proposition~3.10]{GuiUnbounded}, equivalently
\cite[Remark~4.23]{GuiExtensions}.  This gives exactly $e^{itL_0}$
for rotations, not an arbitrarily shifted positive-energy generator.
The normalization is important when using $r^{L_0}$ in annulus
factorizations.
\end{proof}

The construction respects finite orthogonal sums.  This follows either
from the closed-smear identity and uniqueness, or directly from the
diagonal construction.  Thus the theorem applies to all objects of
the grading-restricted unitary category, including finite sums of the
vacuum.  It makes no assertion yet that every simple algebraic module
admits such a positive form.  We establish that remaining assertion
after proving positivity of transport for pairs of unitary modules.
